\documentclass[12pt,a4paper]{article}
\usepackage[utf8]{inputenc}
\usepackage[spanish,es-tabla]{babel}
\usepackage{geometry}
\usepackage{graphicx}
\usepackage{multirow}
\usepackage{tabularx}
\usepackage{mathptmx}
\usepackage{fancyhdr}
\usepackage{array}
\usepackage[table]{xcolor}
\usepackage{titlesec}
\usepackage{setspace}
\usepackage{microtype}
\usepackage{lastpage}
\usepackage{booktabs}
\usepackage{enumitem}
\usepackage{amsmath}
\usepackage{amsfonts}
\usepackage{amssymb}
\usepackage{listings}

\lstset{
  basicstyle=\ttfamily\footnotesize,
  breaklines=true,
  frame=single,
  numbers=none,
  backgroundcolor=\color{gray!5},
  showstringspaces=false
}

\graphicspath{{images/}{../images/}{../../images/}{./}}

\geometry{
    top=4.7cm,
    bottom=2.5cm,
    left=2.5cm,
    right=2.5cm,
    headheight=3.9cm,
    headsep=0.4cm
}

\setstretch{1.15}
\setlength{\parindent}{0pt}
\setlength{\parskip}{0.45em}
\setlength{\tabcolsep}{4pt}

\titleformat{\section}
  {\normalfont\fontsize{12}{14.5}\bfseries}{\thesection.}{0.6em}{\uppercase}
\titlespacing*{\section}
  {0pt}{1.8ex plus 0.4ex minus 0.2ex}{0.6ex plus 0.2ex}

\titleformat{\subsection}
  {\normalfont\fontsize{11.5}{13.5}\bfseries}{\thesubsection.}{0.5em}{}
\titlespacing*{\subsection}
  {0pt}{1.4ex plus 0.3ex minus 0.1ex}{0.4ex plus 0.1ex}

\titleformat{\subsubsection}
  {\normalfont\fontsize{10.5}{12.5}\bfseries\itshape}{\thesubsubsection.}{0.4em}{}
\titlespacing*{\subsubsection}
  {0pt}{1.1ex plus 0.2ex minus 0.1ex}{0.3ex plus 0.1ex}

\newcolumntype{Y}{>{\centering\arraybackslash}X}
\newcolumntype{L}{>{\raggedright\arraybackslash}X}
\newcolumntype{C}{>{\centering\arraybackslash}X}

\definecolor{headerred}{RGB}{180, 0, 0}
\definecolor{darkred}{RGB}{140, 0, 0}
\definecolor{lightgraybg}{RGB}{248, 248, 248}
\definecolor{tableheader}{RGB}{232, 238, 246}

\newcommand{\docCodigo}{TI-SIGE-ENT-S5-001}
\newcommand{\docVersion}{1.0}
\newcommand{\docProceso}{DESARROLLO FRONTEND / MATRIZ DINÁMICA DE ACTIVIDADES}
\newcommand{\docElaboracion}{29/09/2026}
\newcommand{\docRevision}{APROBADO POR LÍDER TÉCNICO}
\newcommand{\docTitulo}{DESARROLLO DE LA MATRIZ DINÁMICA DE ACTIVIDADES (EDT 3.9)}
\newcommand{\docOrganizacion}{FACULTAD DE INGENIERÍA EN SISTEMAS, ELECTRÓNICA E INDUSTRIAL}
\newcommand{\docSubOrganizacion}{CARRERA DE INGENIERÍA DE SOFTWARE}

\newcommand{\encabezadofrecuente}{%
    \renewcommand{\arraystretch}{1.15}%
    \noindent\makebox[\textwidth][c]{%
    \begin{tabularx}{\textwidth}{|>{\centering\arraybackslash}m{2.8cm}|Y|>{\centering\arraybackslash}m{4.8cm}|}
        \hline
        \multirow{4}{2.8cm}{\centering\raisebox{-1.15cm}{\includegraphics[width=2.5cm,height=2.0cm,keepaspectratio]{images/logo.png}}}
        & \textbf{\fontsize{9}{11}\selectfont \docOrganizacion} 
        & \textbf{\fontsize{8}{10}\selectfont CÓDIGO:} \fontsize{8}{10}\selectfont \docCodigo \\
        \cline{3-3}
        & \textbf{\fontsize{8.5}{10.5}\selectfont \docSubOrganizacion} 
        & \textbf{\fontsize{8}{10}\selectfont VERSIÓN:} \fontsize{8}{10}\selectfont \docVersion \\
        \cline{2-3}
        & \textbf{\fontsize{8}{10}\selectfont PROCESO:} \fontsize{8}{10}\selectfont \docProceso 
        & \textbf{\fontsize{8}{10}\selectfont ELABORACIÓN:} \fontsize{8}{10}\selectfont \docElaboracion \\
        \cline{2-3}
        & \textbf{\fontsize{8.5}{10.5}\selectfont \docTitulo} 
        & \textbf{\fontsize{8}{10}\selectfont ÚLTIMA REVISIÓN:} \fontsize{8}{10}\selectfont \textbf{\textcolor{darkred}{\docRevision}} \\
        \hline
    \end{tabularx}%
    }%
}

\pagestyle{fancy}
\fancyhf{}
\fancyhead[C]{\encabezadofrecuente}
\fancyfoot[C]{\fontsize{10}{12}\selectfont Página \thepage\ de \pageref{LastPage}}
\renewcommand{\headrulewidth}{0pt}

\begin{document}

\section{Título}
\textbf{Denominación del Entregable:} Desarrollo e Implementación de la Matriz Dinámica de Actividades y Cronogramas del SGC (EDT 3.9)

\section{Introducción}
La matriz de actividades constituye el núcleo operativo y pedagógico de los formatos institucionales \texttt{UTA-SGC-A-2-1-P7-T1} (Plan de Trabajo) y \texttt{UTA-SGC-A-2-1-P7-T2} (Informe de Cumplimiento). En el proceso tradicional, los docentes elaboraban esta matriz mediante tablas estáticas de Word y hojas de cálculo desvinculadas, lo que generaba errores frecuentes de cálculo en horas, fechas fuera del semestre lectivo y desajustes en el porcentaje de avance.

Para superar estas limitaciones, en el marco de la iteración de la Semana 5 se construyó el componente interactivo de \textbf{Matriz Dinámica de Actividades}. Esta herramienta permite a los docentes agregar, editar y ordenar actividades de forma fluida, seleccionar fechas mediante modales de cronograma inteligentes, asignar medios de verificación y visualizar el recálculo automático de porcentajes y horas asignadas en estricto cumplimiento con la normativa universitaria.

\section{Objetivo}

\subsection{Objetivo General}
Desarrollar e integrar el componente de interfaz y lógica reactiva de la Matriz Dinámica de Actividades en el sistema SIGE, proporcionando cálculo automatizado de avances, validación de cronogramas y sincronización bidireccional entre la planificación docente y los informes de cumplimiento.

\subsection{Objetivos Específicos}
\begin{itemize}[leftmargin=1.5em, itemsep=1.5pt]
    \item Implementar la tabla interactiva de actividades con adición dinámica de filas, eliminación segura y edición en línea sin pérdida de estado.
    \item Desarrollar el modal selector de cronograma (\texttt{CronogramaPickerModal.jsx}) para la asignación ágil de rangos de fechas (Desde -- Hasta) validados con el calendario académico.
    \item Implementar el algoritmo de cálculo instantáneo de horas, ponderaciones acumuladas y avance porcentual global.
    \item Vincular la matriz al catálogo institucional de actividades predefinidas de comisiones, tutorías y gestión académica de la FISEI.
\end{itemize}

\section{Alcance}

\subsection{Alcance Incluido}
\begin{itemize}[leftmargin=1.5em, itemsep=1.5pt]
    \item Componente de matriz dinámica en \texttt{PlanTrabajoView.jsx} e \texttt{InformeView.jsx}.
    \item Selector interactivo de fechas desde/hasta con delimitación del semestre lectivo vigente.
    \item Catálogo de actividades institucionales filtradas por comisión oficial de la FISEI.
    \item Controles de validación de campos obligatorios: actividad, cronograma, recursos y medios de verificación.
\end{itemize}

\subsection{Exclusiones}
\begin{itemize}[leftmargin=1.5em, itemsep=1.5pt]
    \item Carga física de archivos binarios de evidencia (gestionada por el servicio MinIO S3 en EDT 3.11).
    \item Bloqueo de firmas por motor de reglas backend (gestionado en Validation Gates en EDT 3.10).
\end{itemize}

\section{Definiciones}
\begin{itemize}[leftmargin=1.5em, itemsep=2pt]
    \item \textbf{Matriz Dinámica:} Componente de software de tabla reactiva que actualiza estados, cálculos y renderizado sin recargar la página.
    \item \textbf{Ponderación:} Peso porcentual relativo asignado a cada actividad planificada cuya suma total debe alcanzar el 100\%.
    \item \textbf{Medio de Verificación:} Evidencia documental objetiva y auditable que respalda la ejecución de una actividad académica.
    \item \textbf{Cronograma SGC:} Intervalo temporal formal (fecha de inicio y fecha de finalización) en el que se ejecuta cada labor docente.
\end{itemize}

\section{Responsabilidades}
\begin{itemize}[leftmargin=1.5em, itemsep=2pt]
    \item \textbf{Desarrolladora Frontend (Andrea Vásquez):} Maquetación de la matriz, diseño de controles reactivos y vinculación del selector de fechas.
    \item \textbf{Analista UI/UX (Heidi Villavicencio):} Validación de la experiencia de usuario y coherencia con la plantilla impresa oficial A4.
    \item \textbf{Líder de Proyecto / Arquitecto (Steeven Toala):} Revisión del modelo de datos de la matriz y control de calidad del código.
    \item \textbf{Docente Tutor (Ing. Paulo Torres):} Aprobación metodológica y verificación de cumplimiento con el SGC.
\end{itemize}

\section{Normativa Legal y Técnica}
\begin{itemize}[leftmargin=1.5em, itemsep=1.5pt]
    \item \textbf{Reglamento de Régimen Académico Institucional (UTA):} Distribución de carga horaria docente en docencia, investigación y vinculación.
    \item \textbf{Manual de Procesos SGC - FISEI:} Procedimiento P7 para la planificación y seguimiento del desempeño académico.
    \item \textbf{Estándar de Accesibilidad W3C WCAG 2.1:} Accesibilidad en tablas de datos complejas e interactividad por teclado.
\end{itemize}

\section{Método, Políticas y Contenidos}
La implementación se realizó bajo la arquitectura basada en componentes de \textbf{React 19}, aprovechando el estado local y la propagación de eventos para la reactividad en tiempo real:
\begin{enumerate}[leftmargin=1.5em, itemsep=2pt]
    \item \textbf{Gestión de Estado:} Cada fila de la matriz se representa mediante un objeto estructurado que contiene identificador único, descripción de actividad, fechas, responsable institucional, recursos y medios de verificación.
    \item \textbf{Políticas de Validación Temprana:} La matriz impide avanzar a la siguiente sección si existen filas con descripciones en blanco o sin cronograma definido.
    \item \textbf{Integración con Catálogos:} Los docentes pueden autocompletar actividades típicas desde el panel lateral \texttt{ActivityCatalogSidebar.jsx}.
\end{enumerate}

\section{Descripción de actividades}

\subsection{Estructura de Datos de la Matriz Dinámica}
A continuación, se presenta la especificación del esquema de datos que estructura cada elemento de la matriz en el frontend:

\begin{lstlisting}[language=JavaScript, caption={Estructura JSON de una actividad en la matriz dinámica}]
{
  "id": "act-1727618291000",
  "numero": 1,
  "actividad": "Elaboración y actualización de reactivos para exámenes complexivos",
  "desde": "11/08/2026",
  "hasta": "30/09/2026",
  "responsable": "Docente Responsable / Comisión de Titulación",
  "recursos": "Material bibliográfico, plataforma virtual institucional",
  "mediosVerificacion": "Banco de reactivos validado y acta de entrega",
  "metaPlanificada": 100,
  "avanceReal": 100
}
\end{lstlisting}

\subsection{Componentes Desarrollados y Funcionalidades Clave}
\begin{itemize}[leftmargin=1.5em, itemsep=2.5pt]
    \item \textbf{Tabla Dinámica A4:} Diseñada para renderizarse con formato exacto al documento impreso de la UTA, soportando saltos de página fluidos y numeración automática.
    \item \textbf{Selector de Fechas (\texttt{CronogramaPickerModal}):} Permite al docente seleccionar fechas de inicio y fin con un solo clic a partir del calendario semestral oficial preestablecido (agosto -- diciembre 2026).
    \item \textbf{Adición Rápida desde Catálogo:} Mediante el botón \textit{"+ Catálogo"}, el usuario despliega actividades normadas y las inserta con un clic, reduciendo el tiempo de digitación en más de un 70\%.
\end{itemize}

\section{Firmas de Responsabilidad}

\noindent
\renewcommand{\arraystretch}{1.2}
\begin{tabularx}{\textwidth}{|>{\raggedright\arraybackslash}m{2.9cm}|>{\raggedright\arraybackslash}X|>{\centering\arraybackslash}m{4.6cm}|>{\centering\arraybackslash}m{2.3cm}|}
    \hline
    \rowcolor{headerred}
    \textcolor{white}{\textbf{ACCIÓN}} & \textcolor{white}{\textbf{NOMBRE Y CARGO}} & \textcolor{white}{\textbf{FIRMA}} & \textcolor{white}{\textbf{FECHA}} \\
    \hline
    \textbf{Elaborado por:} & Andrea Vásquez\newline \textit{Desarrolladora Frontend / QA} & \parbox[c]{4.4cm}{\centering\vspace{4pt}\includegraphics[height=1.5cm,width=4.0cm,keepaspectratio]{images/firma_heidi.png}\vspace{4pt}} & 29/09/2026 \\
    \hline
    \textbf{Revisado por:} & Steeven Toala\newline \textit{Líder de Proyecto / Arquitecto Software} & \parbox[c]{4.4cm}{\centering\vspace{4pt}\includegraphics[height=1.5cm,width=4.0cm,keepaspectratio]{images/firma_steveen.png}\vspace{4pt}} & 29/09/2026 \\
    \hline
    \textbf{Aprobado por:} & Ing. Paulo Torres\newline \textit{Docente Tutor - Gestión Proyectos Software} & \parbox[c]{4.4cm}{\centering\vspace{10pt}\textbf{\textcolor{gray!80}{\scriptsize [ PENDIENTE DE APROBACIÓN ]}}\vspace{10pt}} & \textbf{\textcolor{gray!90}{Pendiente}} \\
    \hline
\end{tabularx}

\vspace{0.4cm}

\section{Control de Cambios}

\noindent
\renewcommand{\arraystretch}{1.2}
\begin{tabularx}{\textwidth}{|>{\hsize=0.45\hsize\centering\arraybackslash}X|>{\hsize=1.55\hsize\raggedright\arraybackslash}X|>{\hsize=1.2\hsize\raggedright\arraybackslash}X|>{\hsize=0.8\hsize\centering\arraybackslash}X|}
    \hline
    \rowcolor{gray!20}
    \textbf{Versión} & \textbf{Descripción del cambio} & \textbf{Responsable del cambio} & \textbf{Fecha de aprobación} \\
    \hline
    1.0 & Creación e integración del componente Matriz Dinámica de Actividades y selector de cronogramas. & Andrea Vásquez / Steeven Toala & 29/09/2026 \\
    \hline
\end{tabularx}

\vspace{0.4cm}

\section{Anexos}
\subsection*{Anexo 1: Captura del Componente Matriz de Actividades}
\begin{itemize}[leftmargin=1.5em, itemsep=1.5pt]
    \item Archivo fuente: \texttt{frontend/src/components/docente/planificacion/PlanTrabajoView.jsx}
    \item Selector de fechas: \texttt{frontend/src/components/docente/planificacion/CronogramaPickerModal.jsx}
\end{itemize}

\end{document}
